% GDZ PPN236018647,native589middle;HKUST590.
\NumberedItem{117}
\noindent\hspace*{\parindent}Iisdem diebus Pascha Iudaeorum per methodum novam determinare docuimus.
\par\noindent\hfill [1801] (Apr. 1.)
\par\medskip\begingroup\small
\renewcommand{\thefootnote}{**)}
\noindent\hspace*{\parindent}Über die Einrichtung des jüdischen Kalenders dürfte \textsc{Gauss} sich aus \textsc{Christian Wolfs} \textit{Elementa matheseos universae}\footnote{Tomus IV., qui geographiam cum hydrographia, chronologiam etc. complectitur, handelt im Cap. VII. der \textit{Elementa chronologiae}, S. 182, de calendariis iudaico et muhamedano. In dem von \textsc{Gauss} 1800 erworbenen Abdruck des \textsc{Wolf}schen Werkes findet sich eine handschriftliche Eintragung \textit{Die Berechnung des Neumonds Tisri für jedes jüdische Jahr A}, die Werke XI1 zum Abdruck gelangt.} unterrichtet haben. Seine \textit{Berechnung des jüdischen Osterfestes} ist in \textsc{v. Zachs} Monatlicher Correspondenz der Erd- und Himmelskunde \textbf{5}, Mai 1802, S. 435 veröffentlicht worden (Werke VI, S. 80); den ersten Beweis der \textsc{Gauss}schen Regel gab \textsc{Cisa Gresy}\renewcommand{\thefootnote}{***)}\footnote{\textsc{Chevalier Cisa Gresy}, \textit{Démonstration des formules de M. Gauss pour déterminer le jour de Pâque des juifs}, Correspondance astronom. \textbf{1}, 1818, S. 556.} auf Veranlassung des Freiherrn \textsc{von Zach}\renewcommand{\thefootnote}{\ensuremath{\dagger})}\footnote{Siehe ebenda, S. 568.}; \textsc{Ch. Z. Slonimsky}\renewcommand{\thefootnote}{\ensuremath{\dagger\dagger})}\footnote{\textsc{Ch. Z. Slonimsky}, \textit{Eine allgemeine Formel für die gesamte jüdische Kalenderberechnung}, \textsc{Crelles} Journal für Mathematik \textbf{28}, 1844, S. 179.} hat die \textsc{Gauss}sche Formel so ausgestaltet, daß sie auf alle Fragen des jüdischen Kalenders eines Jahres Auskunft erteilt; einen durchsichtigen Beweis gab \textsc{M. Hamburger}\renewcommand{\thefootnote}{\ensuremath{\dagger\dagger\dagger})}\footnote{\textsc{M. Hamburger}, \textit{Ableitung der Gaussschen Formel zur Bestimmung des jüdischen Osterfestes}, ebenda \textbf{116}, 1896, S. 90.}.
\par\noindent\hfill\textsc{Klein.\quad Loewy.}
\par\endgroup
